\section{Conclusion}

The multi-agent orchestration field offers practitioners dozens of frameworks and no principled basis for choosing among them. We addressed that gap through three contributions: (1)~the D-I-T characterization that turns topology selection from guesswork into geometry, (2)~OrchestraBench, a controlled comparison producing 794 runs across three seeds showing that hierarchical reaches 75\% on high-D tasks while debate reaches 77\% on high-I tasks, and (3)~a DIT routing classifier achieving 90\% accuracy (bootstrap 95\% CI: [84\%, 95\%]) against a 33\% random baseline.

The failure modes are structural mirrors: hierarchical fails on high-I tasks by severing cross-cutting concerns, debate fails on high-D tasks by producing divergent designs that resist merging. Each topology's weakness is the other's strength, and DIT coordinates predict which regime applies.

Orchestration topology is not a detail to be decided by convention. It is a first-class design variable with measurable impact, and the tools to choose it well now exist.
